\section{引言}

\subsection{应用背景与问题}

开口或半开口液体运输是实验室自动化、服务机器人和移动操作系统中的基础任务之一。与普通货物运输不同，液体在容器内会对机器人运动产生具有历史记忆的动态响应：平动加速度、转向运动和速度变化会激发自由液面振荡，进而带来溢出风险、运输效率下降和安全隐患。因此，移动机器人在运输液体时不能只追求几何路径跟踪或表面上的轨迹平滑，还需要规划一段与液体动态响应相容的运动。

本文关注的问题可以概括为：在给定安全参考路径的条件下，移动底盘如何在线生成满足底盘运动约束的局部轨迹，同时预测并抑制由机器人运动诱发的液体晃动。该问题处在液体防晃和移动机器人局部规划的交界处，但它既不是通用流体力学建模问题，也不是普通局部规划器的平滑性调参问题。

\subsection{现有路线与缺口}

液体防晃已经在多个平台上被研究。机械臂、SCARA 和操作机器人中的防晃方法表明，显式考虑液体模型和容器运动对液面响应的影响是有价值的；液罐车、车辆和船舶等载液平台研究则说明，液体晃动会影响移动载体的稳定性和安全性。同时，跨平台基础方法如低阶液体模型、输入整形、速度剖面和液面测量为防晃规划提供了模型和评价基础。

然而，对于移动底盘或服务机器人运输开口液体这一场景，已有近邻方法常集中在固定路径或曲线路径设计、给定路径下的速度剖面、离线防晃轨迹优化、给定轨迹后的防晃跟踪控制或特殊主动抑振机构。它们能够从不同角度减少晃动，但往往不是移动机器人导航栈中的在线局部轨迹规划问题。另一方面，普通移动机器人局部规划器能够实时生成可行、平滑甚至避障的轨迹，却通常不传播所运输液体的内部动态状态，因此无法区分“看起来平滑”的运动和“与液体动态记忆相容”的运动。

由此形成本文的核心缺口：移动底盘液体运输需要一种在线滚动的防晃局部轨迹规划方法，在规划阶段显式预测液体状态，并在路径进度、底盘控制和液体响应之间进行权衡。

\subsection{本文方法概述}

为解决上述问题，本文提出晃动感知模型预测轮廓控制方法（Slosh-aware Model Predictive Contouring Control, \SPMPC）。该方法将移动底盘的局部轨迹规划表述为带路径进度变量的 \MPCC{} 问题，并在机器人状态中加入低阶液体模态状态。具体而言，规划器在滚动时域内同时优化路径进度、底盘加速度、角加速度和虚拟路径速度，并利用模态动力学预测纵向和横向激励下的液体响应。

与仅限制速度、加速度或 jerk 的平滑规划不同，\SPMPC{} 将液体模态坐标及其速度作为规划状态，使当前液体相位和历史激励能够影响未来局部运动的选择。与给定轨迹后的防晃跟踪控制不同，\SPMPC{} 在局部轨迹生成阶段就将液体响应纳入目标函数和约束。与离线防晃轨迹优化不同，\SPMPC{} 面向移动底盘在线执行，每个控制周期重新求解有限时域优化问题。

\subsection{贡献与范围}

本文第一版的贡献计划限定为以下三点：

\begin{enumerate}
  \item \textbf{问题与方法表述：} 将移动底盘开口液体运输建模为在线防晃局部轨迹规划问题，提出包含低阶液体模态状态的 alpha-state \MPCC{} 表述。
  \item \textbf{系统实现：} 基于 ROS 和 acados 实现可实时运行的滚动时域局部规划器，在给定安全参考路径上联合优化路径进度、底盘控制和预测液体响应。
  \item \textbf{证据链设计：} 通过 slosh-aware 与 smooth-only 消融、普通局部规划器 baseline 对比以及外部视觉液面评价，验证显式液体状态预测相对于单纯平滑运动的独立价值。
\end{enumerate}

本文不把完整障碍物感知 \MPCC{}、homotopy/corridor 约束、随机机会约束、CBF 主方法、高保真 CFD/FEM/SPH 在线模型或完整液面状态观测器作为第一版贡献。本文的范围是沿给定安全实验室参考路径的在线防晃局部轨迹规划与控制。
